\section{Tutorial 2: Gradient Descent Exercises}

\subsection{Objectives}

This tutorial develops intuition and technical skill for gradient-based optimization:
\begin{itemize}
    \item Compute gradients for common losses
    \item Analyze convergence of gradient descent
    \item Understand the role of step size and conditioning
    \item Connect algebraic updates with geometric intuition
\end{itemize}

\medskip

\noindent
\textbf{Focus.}  
Work through simple but revealing cases where everything can be computed explicitly.

\subsection{Exercise 1: Gradient of a Quadratic Loss}

Let
\[
L(w) = \frac{1}{2}\|Xw - y\|^2, \quad X \in \mathbb{R}^{n \times d}, \; y \in \mathbb{R}^n.
\]

\medskip

\noindent
\textbf{Tasks.}
\begin{enumerate}
    \item Compute $\nabla_w L(w)$.
    \item Write the gradient descent update.
\end{enumerate}

\medskip

\noindent
\textbf{Solution.}

\begin{enumerate}
\item Using matrix calculus:
\[
\nabla_w L(w) = X^\top (Xw - y).
\]

\item Gradient descent:
\[
w_{t+1} = w_t - \eta X^\top (Xw_t - y).
\]
\end{enumerate}

\medskip

\noindent
\textbf{Key takeaway.}  
Quadratic objectives lead to linear gradient updates.

\subsection{Exercise 2: One-Dimensional Gradient Descent}

Consider
\[
L(\theta) = \frac{1}{2} a \theta^2, \quad a > 0.
\]

\medskip

\noindent
\textbf{Tasks.}
\begin{enumerate}
    \item Compute the gradient.
    \item Write the update rule.
    \item Solve the recursion for $\theta_t$.
    \item Determine when the method converges.
\end{enumerate}

\medskip

\noindent
\textbf{Solution.}

\begin{enumerate}
\item 
\[
\nabla L(\theta) = a\theta.
\]

\item 
\[
\theta_{t+1} = \theta_t - \eta a \theta_t = (1 - \eta a)\theta_t.
\]

\item 
\[
\theta_t = (1 - \eta a)^t \theta_0.
\]

\item Convergence requires:
\[
|1 - \eta a| < 1 \quad \Rightarrow \quad 0 < \eta < \frac{2}{a}.
\]
\end{enumerate}

\medskip

\noindent
\textbf{Key takeaway.}  
Step size controls convergence speed and stability.

\subsection{Exercise 3: Effect of Conditioning}

Let
\[
L(w) = \frac{1}{2} w^\top A w,
\]
where $A$ is symmetric positive definite.

\medskip

\noindent
\textbf{Tasks.}
\begin{enumerate}
    \item Compute the gradient.
    \item Explain how eigenvalues of $A$ affect convergence.
\end{enumerate}

\medskip

\noindent
\textbf{Solution.}

\begin{enumerate}
\item 
\[
\nabla L(w) = A w.
\]

\item Let eigenvalues of $A$ be $\lambda_1 \leq \cdots \leq \lambda_d$.

\begin{itemize}
    \item Convergence rate depends on condition number:
    \[
    \kappa = \frac{\lambda_{\max}}{\lambda_{\min}}.
    \]
    \item Large $\kappa$ $\Rightarrow$ slow convergence
\end{itemize}
\end{enumerate}

\medskip

\noindent
\textbf{Geometric interpretation.}  
Level sets are ellipses; gradient descent zig-zags along narrow directions.

\medskip

\noindent
\textbf{Key takeaway.}  
Poor conditioning slows optimization.

\subsection{Exercise 4: Stochastic Gradient Descent}

Consider empirical loss:
\[
L(\theta) = \frac{1}{n} \sum_{i=1}^n \ell_i(\theta).
\]

\medskip

\noindent
\textbf{Tasks.}
\begin{enumerate}
    \item Define the SGD update.
    \item Compare SGD with full gradient descent.
    \item Explain the effect of stochastic noise.
\end{enumerate}

\medskip

\noindent
\textbf{Solution.}

\begin{enumerate}
\item SGD:
\[
\theta_{t+1} = \theta_t - \eta \nabla \ell_i(\theta_t),
\]
where $i$ is sampled randomly.

\item Comparison:
\begin{itemize}
    \item GD: exact gradient, expensive
    \item SGD: noisy gradient, cheap
\end{itemize}

\item Noise effects:
\begin{itemize}
    \item Helps escape saddle points
    \item Introduces fluctuations near minima
\end{itemize}
\end{enumerate}

\medskip

\noindent
\textbf{Key takeaway.}  
SGD trades accuracy for efficiency and exploration.

\subsection{Exercise 5: Gradient Descent as Projection}

Consider linear regression:
\[
\min_w \|Xw - y\|^2.
\]

\medskip

\noindent
\textbf{Tasks.}
\begin{enumerate}
    \item Describe the solution geometrically.
    \item Explain how gradient descent approaches this solution.
\end{enumerate}

\medskip

\noindent
\textbf{Solution.}

\begin{enumerate}
\item The solution is the projection of $y$ onto the column space of $X$.

\item Gradient descent:
\begin{itemize}
    \item Iteratively reduces residual $Xw - y$
    \item Moves toward projection point
\end{itemize}
\end{enumerate}

\medskip

\noindent
\textbf{Key takeaway.}  
Optimization can be interpreted as a geometric projection process.

\subsection{Exercise 6 (Discussion): Why Does GD Work in High Dimensions?}

\textbf{Question.}  
Neural networks are highly non-convex. Why does gradient descent often work well?

\medskip

\noindent
\textbf{Discussion points.}
\begin{itemize}
    \item Many local minima are similarly good
    \item High-dimensional landscapes dominated by saddle points
    \item Overparameterization simplifies optimization
\end{itemize}

\medskip

\noindent
\textbf{Insight.}  
Optimization in modern machine learning behaves differently from classical low-dimensional intuition.

\subsection{Summary}

\begin{itemize}
    \item Gradient descent follows local gradients to minimize loss
    \item Step size controls convergence behavior
    \item Conditioning affects speed of convergence
    \item Stochasticity introduces noise but aids exploration
    \item Optimization has a geometric interpretation
\end{itemize}

\medskip

\noindent
\textbf{Preparation for next lecture.}  
We now study how optimization interacts with generalization and overfitting.